\section*{第 \ref{ch:b1:multivar} 章——二元函数}

\begin{solution}{exo:b1:multivar:1}
$\dfrac{\partial f}{\partial x} = 2xy + y\,\eu^{xy}$,
$\dfrac{\partial f}{\partial y} = x^2 + x\,\eu^{xy}$.

$\dfrac{\partial g}{\partial x} = \dfrac{2x}{x^2+y^2}$,
$\dfrac{\partial g}{\partial y} = \dfrac{2y}{x^2+y^2}$.

$\dfrac{\partial h}{\partial x} = \dfrac{-y/x^2}{1 + y^2/x^2} =
\dfrac{-y}{x^2+y^2}$,
$\dfrac{\partial h}{\partial y} = \dfrac{x}{x^2+y^2}$.
\end{solution}

\begin{solution}{exo:b1:multivar:2}
在极\omterm{prop:b1:vspaces:coordinates}{坐标}下（$\rho \to 0$）：

$\abs{f} = \dfrac{\rho^3\abs{\cos^2\theta\sin\theta}}{\rho^2} \leq
\rho \to 0$：\omterm{def:b1:continuity:continuous}{连续}。

$g = \cos\theta\sin\theta$：与 $\rho$ 无关，
沿不同射线取不同值（参见\
\cref{ex:b1:multivar:trap}）：极限不存在，不\omterm{def:b1:continuity:continuous}{连续}。

$\abs h \leq \dfrac{\rho^3(\abs{\cos^3\theta} +
\abs{\sin^3\theta})}{\rho^2} \leq 2\rho \to 0$：\omterm{def:b1:continuity:continuous}{连续}。
\end{solution}

\begin{solution}{exo:b1:multivar:3}
沿 $y = mx$：$f(x, mx) = \dfrac{m x^3}{x^4 + m^2 x^2} =
\dfrac{mx}{x^2 + m^2} \to 0$（$m \neq 0$；沿 $y = 0$ 与
$y$ 轴时 $f = 0$）。故所有直线上的极限均为 $0$。但在
抛物线 $y = x^2$ 上：
\[
f(x, x^2) = \frac{x^4}{x^4 + x^4} = \frac12 :
\]
点列 $\bigl(\frac1n, \frac{1}{n^2}\bigr) \to (0,0)$ 使得
$f \to \frac12 \neq 0$。因此不\omterm{def:b1:continuity:continuous}{连续}：仅沿直线
检验不足以判定二元函数的极限。
\end{solution}

\begin{solution}{exo:b1:multivar:4}
$\frac{\partial f}{\partial x} = 3x^2y^2 + y\cos(xy)$；于是
\[
\frac{\partial^2 f}{\partial y\,\partial x}
= 6x^2 y + \cos(xy) - xy\sin(xy) .
\]
$\frac{\partial f}{\partial y} = 2x^3 y + x\cos(xy)$；于是
\[
\frac{\partial^2 f}{\partial x\,\partial y}
= 6x^2 y + \cos(xy) - xy\sin(xy) :
\]
二者相等，正如施瓦茨定理所保证的。
\end{solution}

\begin{solution}{exo:b1:multivar:5}
由\cref{thm:b1:multivar:chain}，取 $(x, y) = (\cos t, \sin t)$：
\[
g'(t) = -\sin t\,\frac{\partial f}{\partial x}(\cos t, \sin t)
+ \cos t\,\frac{\partial f}{\partial y}(\cos t, \sin t)
= \bigl\langle \nabla f,\ (-\sin t, \cos t)\bigr\rangle .
\]
在 $g$ 的极值点处 $g'(t) = 0$，故 $\nabla f$ 与
圆的切向量 $(-\sin t, \cos t)$ 正交，从而与
半径向量 $(\cos t, \sin t)$ 平行（平面上，单位向量的
\omterm{def:b1:euclid:orthogonal}{正交补}是其所张成的直线，取其垂直方向）。
这是拉格朗日乘子法的
最简单的实例。
\end{solution}

\begin{solution}{exo:b1:multivar:6}
对 $f$：$\nabla f = (2x + y - 3,\; x + 2y) = 0$，得 $y = -\frac x2$，且
$2x - \frac x2 = 3$，得 $x = 2$，$y = -1$。二阶：$r = 2$，$s =
1$，$t = 2$：$rt - s^2 = 3 > 0$，$r > 0$：局部（事实上是整体的——
因为是二次函数）极小值，在 $(2, -1)$ 处取值 $f(2,-1) = -3$。

对 $g$：$\nabla g = (2x,\; -2y + 4) = 0$：\omterm{thm:b1:multivar:critical}{临界点} $(0, 2)$；$r = 2$，$s
= 0$，$t = -2$：$rt - s^2 = -4 < 0$：\omterm{met:b1:multivar:monge}{鞍点}。

对 $h$：$\nabla h = \bigl(4x^3 - 4(x - y),\; 4y^3 + 4(x-y)\bigr) = 0$。
两式相加得 $x^3 + y^3 = 0$，故 $y = -x$；代入得 $4x^3 - 8x = 0$：
$x \in \{0, \pm\sqrt2\}$。\omterm{thm:b1:multivar:critical}{临界点}为 $(0,0)$，
$(\sqrt2, -\sqrt2)$、$(-\sqrt2, \sqrt2)$。
在 $(\pm\sqrt2, \mp\sqrt2)$ 处：$r = 12x^2 - 4 = 20$，$s = 4$，$t =
20$，$rt - s^2 > 0$，$r > 0$：局部极小值（取值 $h = 4 + 4 - 16 =
-8$）。在 $(0,0)$ 处：$r = t = -4$，$s = 4$：$rt - s^2 = 0$，无法判定。
直接考察 $h(x, x) = 2x^4 > 0$ 与 $h(x, -x) = 2x^4 - 8x^2 < 0$，
当 $x \neq 0$ 充分小时，每个\omterm{def:b1:topology:open}{邻域}内函数值正负皆有——
故原点处无极值。
\end{solution}

\begin{solution}{exo:b1:multivar:7}
$x^2 - y^2$：其二阶\omterm{def:b1:multivar:partial}{偏导数}分别为 $2$ 与 $-2$，和为 $0$。$xy$：两个纯
二阶\omterm{def:b1:multivar:partial}{偏导数}皆为零。$\eu^x\cos y$：$\frac{\partial^2}{\partial
x^2} = \eu^x\cos y$, $\frac{\partial^2}{\partial y^2} =
-\eu^x\cos y$：和为 $0$。$\ln(x^2+y^2)$：由
\cref{exo:b1:multivar:1}，
\[
\frac{\partial^2}{\partial x^2}\ln(x^2+y^2)
= \frac{2(x^2+y^2) - 4x^2}{(x^2+y^2)^2}
= \frac{2(y^2 - x^2)}{(x^2+y^2)^2},
\]
而 $y$ 的情形与之相反，和为 $0$。
\end{solution}

\begin{solution}{exo:b1:multivar:8}
\emph{\omterm{def:b1:linmaps:projection}{投影}：}平面 $P$ 的法向量为 $n = (1, 2, -1)$，且
过点 $A = (4, 0, 0)$。距原点最近的点是
$O$ 的\omterm{def:b1:linmaps:projection}{投影}：$p = O + \frac{\langle A - O, n\rangle}{\norm
n^2}\,n = \frac{4}{6}(1,2,-1) = \bigl(\frac23, \frac43,
-\frac23\bigr)$，距离为 $\frac{4}{\sqrt 6}$。（公式：
原点到平面 $\langle X, n \rangle = c$ 的距离为
$\frac{\abs c}{\norm n}$，其中 $c = 4$。）

\emph{极小化：}$f(x,y) = x^2 + y^2 + (x + 2y - 4)^2$。记
$w = x + 2y - 4$ 为最后一个因子：
\[
\nabla f = \bigl(2x + 2w,\; 2y + 4w\bigr) = 0
\iff x = -w \text{ 且 } y = -2w .
\]
代入 $w$ 的定义：$w = -w - 4w - 4$，于是 $w =
-\frac23$，从而 $x = \frac23$，$y = \frac43$，且 $z = w =
-\frac23$。与\omterm{def:b1:linmaps:projection}{投影}所得为同一点，距离为
$\sqrt{\frac49 + \frac{16}{9} + \frac49} = \frac{2\sqrt 6}{3} =
\frac{4}{\sqrt6}$；这是极小值点，因为在无穷远处 $f \to \infty$
（正定二次项与线性项之和）。
\end{solution}

\begin{solution}{exo:b1:multivar:9}
记 $\rho = \sqrt{x^2+y^2}$，有 $\frac{\partial \rho}{\partial x} =
\frac{x}{\rho}$，于是
\[
\frac{\partial f}{\partial x} = \varphi'(\rho)\,\frac{x}{\rho},
\qquad
\frac{\partial^2 f}{\partial x^2}
= \varphi''(\rho)\,\frac{x^2}{\rho^2}
+ \varphi'(\rho)\,\Bigl(\frac{1}{\rho} -
\frac{x^2}{\rho^3}\Bigr),
\]
（由商法则与链式法则）。将关于 $y$ 的对称表达式相加：
$\varphi''$ 项合并得 $\frac{x^2 + y^2}{\rho^2} = 1$，
$\varphi'$ 项合并得 $\frac{2}{\rho} - \frac{x^2 +
y^2}{\rho^3} = \frac{1}{\rho}$：
\[
\Delta f = \varphi'' + \frac{\varphi'}{\rho} .
\]
径向调和函数：求解 $\varphi'' + \frac{\varphi'}{\rho} = 0$，即
一阶方程 $u' + \frac u\rho = 0$，令 $u = \varphi'$，
解得 $u = \frac{c}{\rho}$（见\cref{thm:b1:diffeq:homogeneous1}），
再\omterm{thm:b1:integration:def}{积分}得 $\varphi = c\ln\rho + d$。径向调和函数为
$c\,\ln\sqrt{x^2 + y^2} + d$——即
\cref{exo:b1:multivar:7} 中的对数位势与常数，别无其他。
\end{solution}

\begin{solution}{exo:b1:multivar:10}
对 $z = xy$ 在 $(1,1)$ 处：\omterm{def:b1:multivar:partial}{偏导数}为 $y = 1$ 与 $x = 1$，切
平面为 $z = 1 + (x - 1) + (y - 1) = x + y - 1$。水平
切平面意味着两个\omterm{def:b1:multivar:partial}{偏导数}同时为零，即 $\nabla f = 0$：
由\cref{ex:b1:multivar:extrema}可知，恰为\omterm{thm:b1:multivar:critical}{临界点}
$(0, 0)$ 与 $(1, 1)$，对应水平切平面 $z = 0$ 与 $z =
-1$。“水平切平面”与“\omterm{thm:b1:multivar:critical}{临界点}”
是同一概念，分别在图形与公式中呈现。
\end{solution}

\begin{solution}{exo:b1:multivar:11}
\begin{enumerate}
  \item 固定 $y$，一元函数 $x \mapsto
        f(x,y)$ 在 $\R$ 上导数恒为零，从而为常值函数
        （见\cref{thm:b1:derivative:mvt}）：$f(x, y) = f(0, y)$ 对
        所有 $x$：$f$ 只依赖于 $y$。
  \item 设 $g(u, v) = f\bigl(\frac{u+v}2, \frac{u-v}2\bigr)$。
        由链式法则（见\cref{thm:b1:multivar:chain}，
        对变量 $v$ 应用，视 $u$ 固定）：
        \[
        \frac{\partial g}{\partial v}
        = \frac12\,\frac{\partial f}{\partial x}
        - \frac12\,\frac{\partial f}{\partial y} = 0 .
        \]
        由 (1)，$g$ 只依赖于 $u$：$g(u, v) = \varphi(u)$
        其中 $\varphi$ 为 $C^1$ 类函数，并逆用
        变量代换（$u = x + y$，$v = x - y$），
        \[
        f(x, y) = \varphi(x + y).
        \]
        反之，每个这样的 $f$ 均满足 $f_x = f_y =
        \varphi'$：所有解恰为 $C^1$ 类函数，
        其自变量为 $x + y$。
\end{enumerate}
\end{solution}

\begin{solution}{exo:b1:multivar:12}
在开圆盘内，极值点必为\omterm{thm:b1:multivar:critical}{临界点}：$\nabla f = (y,
x) = 0$ 仅在 $(0,0)$ 处成立，此处 $f = 0$；它是\omterm{met:b1:multivar:monge}{鞍点}（$f(\pm
\varepsilon, \pm\varepsilon) = \varepsilon^2 > 0 > -\varepsilon^2
= f(\pm\varepsilon, \mp\varepsilon)$），故该点不取极值。由
在本阶段不加证明地承认的最值定理可知，上下界均可取到，
且必在\omterm{def:b1:topology:closure}{边界}圆周上取到。在\omterm{def:b1:topology:closure}{边界}上，借助
\cref{exo:b1:multivar:5}，
\[
f(\cos t, \sin t) = \cos t\sin t = \frac{\sin 2t}{2}
\in \intcc{-\tfrac12}{\tfrac12},
\]
最大值 $\frac12$ 在 $t = \frac\pi4, \frac{5\pi}4$ 处取得（即点
$\pm\frac{1}{\sqrt2}(1,1)$），最小值 $-\frac12$ 在 $t =
\frac{3\pi}4, \frac{7\pi}4$ 处取得（即点 $\pm\frac{1}{\sqrt2}(1,-1)$）。
全局最大值为 $\frac12$，全局最小值为 $-\frac12$。
\end{solution}

\begin{solution}{pb:b1:multivar:1}
\textbf{1.} 展开右边：
\[
r\Bigl(h + \frac srk\Bigr)^{\!2} + \frac{rt - s^2}{r}k^2
= rh^2 + 2shk + \frac{s^2}{r}k^2 + \frac{rt - s^2}{r}k^2
= q(h,k) .
\]
若 $rt - s^2 > 0$：两个平方项都带有 $r$ 的符号因子，
且 $q(h,k) = 0$ 迫使 $k = 0$，进而 $h = 0$，故
除原点外严格取 $r$ 的符号。若 $rt - s^2 < 0$：$q(1, 0) = r$，而
$q\bigl(-\frac sr, 1\bigr) = \frac{rt - s^2}{r}$ 的符号与之
相反：两种符号均会出现。

\textbf{2.} 若 $r = 0 \neq t$，交换 $h$ 与 $k$ 的角色
（对 $k$ 配方），结论相同；注意
$rt - s^2 = -s^2 < 0$ 只要 $s \neq 0$ 即成立，且 $q =
2shk + tk^2$ 此时取两种符号（$k$ 相对 $h$ 很小）。若
$r = t = 0$：$q = 2shk$，两种符号均出现，当且仅当 $s \neq 0$，亦即
$rt - s^2 = -s^2 < 0$。小结：$q$ 取两种符号 $\iff rt -
s^2 < 0$；$q$ 仅在原点处为零（定号）$\iff rt -
s^2 > 0$（这蕴含 $r \neq 0$，因 $r = 0$ 时 $q(1,0) =
0$）。

\textbf{3.} 设 $f_x = rx + sy + \beta$ 与 $f_y = sx + ty +
\gamma$，直接展开可得
\[
f(a + h, b + k) = f(a, b) + h\,f_x(a,b) + k\,f_y(a,b)
+ \tfrac12 q(h, k),
\]
不含更高次项（该函数是次数为
$2$ 的\omterm{def:b1:poly:def}{多项式}）。在\omterm{thm:b1:multivar:critical}{临界点}处线性项为零：$f(X_0 + H)
- f(X_0) = \frac12 q(H)$ \emph{精确成立}，因此第 1--2 问的
符号分析给出如下分类：严格整体极小（$rt - s^2 >
0$，$r > 0$），严格整体极大（$rt - s^2 > 0$，$r < 0$），
\omterm{met:b1:multivar:monge}{鞍点}——任一\omterm{def:b1:topology:open}{邻域}内两种符号均出现——当 $rt - s^2 <
0$ 时。此即对二次函数证明了\cref{met:b1:multivar:monge}
所述的结论。

\textbf{4.} 二次型 $q - c(h^2 + k^2)$ 的系数为 $r - c$，$s$，
$t - c$。取 $c = \frac{rt - s^2}{r + t}$（注意 $t > 0$，因为
$rt > s^2 \geq 0$ 且 $r > 0$）：
\[
(r - c)(t - c) - s^2 = rt - s^2 - c(r + t) + c^2 = c^2 \geq 0,
\]
且 $r - c \geq 0$（$c \leq r \iff rt - s^2 \leq r^2 + rt$，
成立）。若 $r - c > 0$，由第 1 问的配方可知
$q - c(h^2 + k^2) \geq 0$；若 $r - c = 0$，则 $s = 0$（由
$-s^2 = -c^2 \leq \dots$ 迫使所列量
$\geq 0$，且第一个因子为 $0$）。此时二次型为 $(t - c)k^2 \geq
0$。两种情形下均有 $q(h,k) \geq c(h^2 + k^2)$。

\textbf{5.} 由第 3--4 问，$f(X) = f(X_0) + \frac12 q(X -
X_0) \geq f(X_0) + \frac c2\norm{X - X_0}^2 \to +\infty$ 当
$\norm X \to \infty$ 时 $f$ 为强制函数，且该不等式
对 $X \neq X_0$ 严格成立：\omterm{thm:b1:multivar:critical}{临界点}即唯一的整体
最小点。至于 $f = x^3 + y^3 - 3xy$ 这一三次函数（
见\cref{ex:b1:multivar:extrema}），上述结论不再成立：$f(x,0)
= x^3 \to -\infty$，且 $(1,1)$ 处的局部极小值
并非整体的——失效的正是二次展开的精确性。

\textbf{6.} 展开平方范数：
\[
f(u,v) = u^2\norm{C_1}^2 + 2uv\,\langle C_1, C_2\rangle +
v^2\norm{C_2}^2 - 2u\,\langle C_1, b\rangle - 2v\,\langle C_2,
b\rangle + \norm b^2 ,
\]
这是关于 $(u, v)$ 的二次函数，其二次部分为
$\frac12(ru^2 + 2suv + tv^2)$，其中 $r = 2\norm{C_1}^2$，$s =
2\langle C_1, C_2\rangle$，$t = 2\norm{C_2}^2$。

\textbf{7.} $rt - s^2 = 4\bigl(\norm{C_1}^2\norm{C_2}^2 -
\langle C_1, C_2\rangle^2\bigr) \geq 0$，由柯西–施瓦茨不等式
（见\cref{thm:b1:euclid:cs}），等号成立当且仅当 $C_1, C_2$
成比例（或其中一个为零），亦即该对向量
线性相关。因此 $rt - s^2 > 0 \iff (C_1, C_2)$ \omterm{def:b1:vspaces:free}{线性无关}。

\textbf{8.} $\nabla f = 0$ 写为
\[
\norm{C_1}^2 u + \langle C_1, C_2\rangle v = \langle C_1,
b\rangle,
\qquad
\langle C_1, C_2\rangle u + \norm{C_2}^2 v = \langle C_2,
b\rangle,
\]
即以 格拉姆矩阵为系数矩阵的\omterm{pb:b1:multivar:1}{正规方程}，它们表明
$\langle uC_1 + vC_2 - b,\ C_i\rangle = 0$ 对 $i = 1, 2$ 成立，亦即
$b - p \perp \operatorname{Vect}(C_1, C_2)$，其中 $p = uC_1 +
vC_2$。由第一部分（第 3、5 问）知唯一的\omterm{thm:b1:multivar:critical}{临界点}即
唯一的整体最小点。由
\cref{thm:b1:euclid:projection}，“$p \in
F$，$b - p \perp F$”刻画出 $p$ 为
$b$ 在 $F = \operatorname{Vect}(C_1, C_2)$ 上的\omterm{thm:b1:euclid:projection}{正交投影}。

\textbf{9.} $\norm{b - p}^2 = \norm b^2 - 2\langle b, p\rangle +
\norm p^2$，由 $\langle p, b - p\rangle = 0$ 可得 $\norm p^2 =
\langle p, b\rangle$，故极小值为 $\norm b^2 - \langle p,
b\rangle$。由毕达哥拉斯定理：$\norm b^2 = \norm p^2 + \norm{b - p}^2$
——数据向量分解为其可解释部分 $p$ 与
残差部分 $b - p$，二者互相正交。

\textbf{10.} $E(\alpha, \beta) = \norm{\alpha C_1 + \beta C_2 -
b}^2$，取前述的 $C_1, C_2, b$；
第 8 问的\omterm{pb:b1:multivar:1}{正规方程}逐分量写为
\[
n\alpha + \Bigl(\sum x_i\Bigr)\beta = \sum y_i,
\qquad
\Bigl(\sum x_i\Bigr)\alpha + \Bigl(\sum x_i^2\Bigr)\beta = \sum
x_i y_i .
\]

\textbf{11.} 两边除以 $n$：$\alpha + \beta\bar x = \bar y$
及 $\alpha\bar x + \beta\bigl(v_x + \bar x^2\bigr) = c_{xy} +
\bar x\bar y$。将 $\alpha = \bar y - \beta\bar x$ 代入
第二式得 $\beta v_x = c_{xy}$，于是
\[
\beta = \frac{c_{xy}}{v_x}, \qquad \alpha = \bar y - \beta\bar
x ,
\]
而第一式恰好表明 $(\bar x, \bar y)$
位于该直线上。

\textbf{12.} $v_x = \frac1n\sum(x_i - \bar x)^2 \geq 0$，其为零
当且仅当每个 $x_i$ 都等于 $\bar x$。又 $rt - s^2 = 4\bigl(n\sum
x_i^2 - (\sum x_i)^2\bigr) = 4n^2 v_x$：第 7 问的\omterm{def:b1:vspaces:free}{线性无关}条件
即 $v_x > 0$，亦即 $x_i$ 不全相等。

\textbf{13.} 取 $\alpha = \bar y - \beta\bar x$，将
数据中心化（$\tilde x_i = x_i - \bar x$，$\tilde y_i = y_i - \bar y$）：
\[
E(\alpha, \beta) = \sum_i(\tilde y_i - \beta\tilde x_i)^2
= n\,v_y - 2\beta\,n\,c_{xy} + \beta^2 n\,v_x ,
\]
在 $\beta = c_{xy}/v_x$ 处取最小值，最小值为 $E_{\min} =
n\bigl(v_y - \frac{c_{xy}^2}{v_x}\bigr) = n v_y(1 - \rho^2)$。
由 $E_{\min} \geq 0$ 及 $v_y > 0$ 得 $\rho^2 \leq 1$，且
$\abs\rho = 1$ 当且仅当 $E_{\min} = 0$，亦即每个残差
为零：各点恰好位于该直线上。

\textbf{14.} $n = 4$: $\bar x = \frac64 = 1.5$, $\bar y = 2$,
$\sum x_i^2 = 14$，故 $v_x = 3.5 - 2.25 = 1.25$，$\sum x_iy_i =
0 + 1 + 6 + 12 = 19$，故 $c_{xy} = 4.75 - 3 = 1.75$。因此
\[
\beta = \frac{1.75}{1.25} = 1.4,
\qquad
\alpha = 2 - 1.4\times1.5 = -0.1 :
\qquad y = 1.4\,x - 0.1 .
\]
拟合值 $-0.1,\ 1.3,\ 2.7,\ 4.1$；残差 $0.1,\ -0.3,\
0.3,\ -0.1$；$E_{\min} = 0.01 + 0.09 + 0.09 + 0.01 = 0.2$
（验算：$v_y = \frac{26}4 - 4 = 2.5$ 及 $4(2.5 -
\frac{1.75^2}{1.25}) = 4(2.5 - 2.45) = 0.2$）。

\textbf{15.} 第 10 题的两个\omterm{pb:b1:multivar:1}{正规方程}恰为
$\sum_i\varepsilon_i = 0$ 及 $\sum_i x_i\varepsilon_i =
0$：残差向量与 $C_1 = (1, \dots, 1)$ 正交，
也与 $C_2 = (x_i)$ 正交——即与整个模型空间正交。特别地，
最优直线总是使其误差相互平衡：误差之和
为零。

\textbf{16.} $\frac{\dd}{\dd\alpha}\sum(y_i - \alpha)^2 =
-2\sum(y_i - \alpha) = 0$ 得 $\alpha = \bar y$（且
二阶导数 $2n > 0$ 说明它是全局极小值点，
该函数是单变量的强制二次函数）。极小值为
$\sum(y_i - \bar y)^2 = n v_y$，方差即为常数模型的
不可约二次误差。

\textbf{17.} $\frac{\dd}{\dd\beta}\sum(y_i - \beta x_i)^2 =
-2\sum x_i(y_i - \beta x_i) = 0$ 得 $\beta_0 = \frac{\sum
x_iy_i}{\sum x_i^2}$。将两个斜率都用中心化量表示
：$\beta_0 = \frac{c_{xy} + \bar x\bar y}{v_x + \bar
x^2}$ 等于 $\frac{c_{xy}}{v_x}$ 当且仅当 $c_{xy}\bar x^2 = \bar
x\bar y\,v_x$，即 $\bar x = 0$（中心化的横\omterm{prop:b1:vspaces:coordinates}{坐标}）或
$\bar y = \beta\bar x$——后者意味着 $\alpha = 0$：
\omterm{pb:b1:multivar:1}{回归直线}已经过原点。

\textbf{18.} $g(\tau) = \norm{M - P_0}^2 - 2\tau\langle M - P_0,
w\rangle + \tau^2$ 在 $\tau^* = \langle M - P_0,
w\rangle$ 时取最小值，其值为 $d(M,D)^2 = \norm{M - P_0}^2 - \langle M -
P_0, w\rangle^2$。在平面中，将 $w$ 扩充为一组标准正交
基 $(w, n)$，其中 $n = \frac{(a, b)}{\sqrt{a^2+b^2}}$（
即 $D$ 的单位法向量）。于是 $\norm{M - P_0}^2 = \langle M - P_0,
w\rangle^2 + \langle M - P_0, n\rangle^2$，故 $d(M, D) =
\abs{\langle M - P_0, n\rangle}$，并利用 $aP_{0x} + bP_{0y} =
-c$，得
\[
d = \frac{\abs{a x_0 + b y_0 + c}}{\sqrt{a^2 + b^2}} .
\]

\textbf{19.} 由 $E_{\min}/n = v_y - \frac{c_{xy}^2}{v_x}$ 及
$\beta^2 v_x = \frac{c_{xy}^2}{v_x}$，得
\[
v_y = \beta^2 v_x + \frac{E_{\min}}n :
\]
总方差 $=$ 沿拟合直线的方差 $+$ 残差
方差。（两边除以 $v_y$，得 $1 = \rho^2 + (1 - \rho^2)$，即
直线所“解释”的方差份额为 $\rho^2$。）

\textbf{20.} 模型空间为
$\operatorname{Vect}\bigl(C_1, C_2, C_3\bigr)$，其中 $C_1 = (1)$，
$C_2 = (x_i)$，$C_3 = (x_i^2)$；极小化 $\norm{aC_1 + bC_2 +
cC_3 - b}^2$ 时，与第 8 题完全一样，归结为 $3\times3$
格拉姆方程组，其矩阵元素为 $\langle C_i, C_j\rangle =
\sum_k x_k^{\,i+j-2}$，即上面所列的矩量矩阵。它
可逆当且仅当 $(C_1, C_2, C_3)$ \omterm{def:b1:vspaces:free}{线性无关}
（见\cref{exo:b1:euclid:11}）；关系式 $a + bx_k + cx_k^2 = 0$
对所有 $k$ 成立，则每个 $x_k$ 都是同一\omterm{def:b1:poly:def}{多项式}的根，该\omterm{def:b1:poly:def}{多项式}
次数 $\leq 2$，由于三个取值互异，这是不可能的，除非 $a
= b = c = 0$。试比较在\cref{ch:b1:det}中的周末问题的矩量矩阵，
其行列式恰为
范德蒙德值的平方。

\textbf{21.} 第一个：$r = 2, s = 1, t = 2$，$rt - s^2 = 3 > 0$，
$r > 0$：原点处为严格全局极小值点。第二个：$r = 2, s
= 3, t = 2$，$rt - s^2 = -5 < 0$，为\omterm{met:b1:multivar:monge}{鞍点}。第三个：$r = s = t =
2$，$rt - s^2 = 0$：退化——但 $x^2 + 2xy + y^2 = (x +
y)^2 \geq 0$ 在整条直线 $y = -x$ 上为零：这是一个
（非严格）全局极小值，沿一条直线取得，行列式
判别法无法识别它。

\textbf{22.} 新的和式（$n = 5$）：$\bar x = 3.2$，$\bar y =
1.6$，$\sum x_iy_i = 19$，故 $c_{xy} = 3.8 - 5.12 = -1.32$，
$\sum x_i^2 = 114$，故 $v_x = 22.8 - 10.24 = 12.56$。新的斜率为：
\[
\beta = \frac{-1.32}{12.56} \approx -0.105 :
\]
仅一个点就使明显递增的趋势（$\beta = 1.4$）变为
略微递减的趋势。平方误差对一个残差只计为
$\varepsilon^2$，因此单个远离的点——即
$\abs{x_{i} - \bar x}$ 大\emph{且}残差大——同时主导
$c_{xy}$ 与 $v_x$：\omterm{pb:b1:multivar:1}{最小二乘}法是有效的，但
不稳健。

\textbf{23.} 记 $W = \sum w_i$，并定义 $\bar x_w =
\frac1W\sum w_ix_i$，$\bar y_w$ 同理，$v_x^w = \frac1W\sum
w_ix_i^2 - \bar x_w^2$，$c^w_{xy} = \frac1W\sum w_ix_iy_i - \bar
x_w\bar y_w$。\omterm{def:b1:logic:map}{映射} $\langle u, v\rangle_w = \sum_i w_iu_iv_i$
是 $\R^n$ 上的\omterm{def:b1:euclid:def}{内积}（$w_i > 0$ 保证正定性），故
第二部分的结论原样适用，\omterm{pb:b1:multivar:1}{正规方程}两边除以 $W$
后化为
\[
\alpha + \beta\bar x_w = \bar y_w,
\qquad
\alpha\bar x_w + \beta(v^w_x + \bar x_w^2) = c^w_{xy} + \bar
x_w\bar y_w ,
\]
由此得 $\beta = c^w_{xy}/v^w_x$ 及 $\alpha = \bar y_w -
\beta\bar x_w$。公式形式不变，只是所有平均都改为加权平均。

\textbf{24.} $E_{\min} = \sum\varepsilon_i^2 = 0$ 可得
每个 $\varepsilon_i = 0$：所有点恰在直线上。
由第 13 题可知这正是 $\abs\rho = 1$ 的情形。检验：对 $(1,1), (2,3),
(3,5)$：$\bar x = 2$，$\bar y = 3$，$v_x = \frac{14}3 - 4 =
\frac23$，$c_{xy} = \frac{22}3 - 6 = \frac43$：$\beta = 2$，
$\alpha = -1$，且对全部三个点确有 $y_i = 2x_i - 1$；
$v_y = \frac{35}3 - 9 = \frac83$ 及 $\rho^2 =
\frac{(4/3)^2}{(2/3)(8/3)} = 1$。

\textbf{25.} (i) 对二次函数，二阶展开
是一个\emph{恒等式}，因此对二次型符号的研究
——纯代数内容，即第 1--2 题——对\omterm{thm:b1:multivar:critical}{临界点}作出分类，
无需借助在本阶段不加证明地承认的泰勒定理。(ii) \omterm{pb:b1:multivar:1}{正规方程}表明
“残差与模型空间正交”，因此解析
极小值\emph{恰为}数据向量的\omterm{thm:b1:euclid:projection}{正交投影}：
微\omterm{thm:b1:integration:def}{积分}与欧氏几何计算的是同一对象。(iii)
相关系数 $\rho = c_{xy}/\sqrt{v_xv_y}$ 度量了
直线所解释的方差份额 $\rho^2$，并由
柯西–施瓦茨不等式得 $\abs\rho \leq 1$，等号仅在
数据共线时成立。(iv) 基与\omterm{def:b1:vspaces:free}{线性无关}性（第 18 章），
格拉姆矩阵、矩量矩阵及其行列式（第 22--23 章），
以及\omterm{thm:b1:euclid:projection}{正交投影}（第 23 章）均重新出现，作为同一算法的
工作部件。第二至三部分
发展\emph{\omterm{pb:b1:multivar:1}{最小二乘}法}（勒让德、高斯），并借助其
\emph{\omterm{pb:b1:multivar:1}{正规方程}}。
\end{solution}
